\documentclass[12pt]{article}
 
\usepackage[utf8]{inputenc}
\usepackage[margin=1in]{geometry}
\usepackage{amsmath}
\usepackage{amssymb}
\usepackage{graphicx}
\usepackage{booktabs}
\usepackage{hyperref}
\usepackage{natbib}
\usepackage{setspace}
\usepackage{float}
\usepackage{listings}
\usepackage{xcolor}
\usepackage{subcaption}
 
% For better code formatting
\lstset{
    basicstyle=\ttfamily\small,
    breaklines=true,
    breakatwhitespace=true,
    captionpos=b,
    frame=single,
    backgroundcolor=\color{gray!10}
}
 
% Set line spacing
% \onehalfspacing
 
% Title and author
\title{Automated Assessment of Large Language Model Code Suggestions Against NIST Cybersecurity Standards: A Comparative Analysis of Manual and AI-Driven Approaches}
 
\author{Dishita Sharma}
 
\date{\today}
 
\begin{document}
 
\maketitle
 
\begin{abstract}
 
Large Language Models (LLMs) have become increasingly prevalent in software development, with tools like ChatGPT and GitHub Copilot serving as common code assistants. However, the security implications of LLM-generated code remain underexplored. In this project, we present an automated approach to assess whether code suggestions from ChatGPT align with the National Institute of Standards and Technology (NIST) Cybersecurity Framework guidelines. We compare the results of automated assessment using multiple modern LLMs (including Claude, GPT-4, and open-source alternatives via OpenRouter) against manual annotations from previous work, analyzing 338 developer conversations from a public Hugging Face dataset. Our findings aim to evaluate the accuracy and efficiency of AI-driven code security assessment compared to human expert review, while also examining the trade-off between assessment quality and computational cost across different LLM providers.
 
\end{abstract}
 
\section{Introduction}
 
The integration of Large Language Models into software development workflows has fundamentally changed how developers write and validate code. Services like ChatGPT, GitHub Copilot, and similar tools now provide real-time code suggestions to millions of developers daily. While these tools increase productivity, they also introduce potential security risks: LLMs are trained on publicly available code repositories, which inevitably contain security vulnerabilities and poor security practices \citep{pearce2022asleep, sandoval2023empirical}.
 
Understanding whether LLM-generated code aligns with established security and privacy standards is critical for organizations adopting these tools. Two semesters ago, our team conducted a manual analysis of 338 developer-ChatGPT conversations from a public Hugging Face dataset, evaluating them against the National Institute of Standards and Technology (NIST) Cybersecurity Framework \citep{nist2018framework}. This manual approach required extensive human effort to classify conversations, identify relevant security and privacy questions, and manually evaluate responses against established best practices.
 
\subsection{Motivation and Research Questions}
 
The primary motivation for this project is to automate the security assessment process using modern LLMs as evaluators, rather than relying solely on manual human review. This raises several important questions:
 
\begin{enumerate}
    \item Can modern LLMs accurately replicate the manual assessment process for evaluating code security?
    \item What are the trade-offs between using expensive, high-capability models (e.g., Claude) versus cheaper, open-source alternatives?
    \item How consistent are LLM-based assessments with expert human judgment?
\end{enumerate}
 
Our hypothesis is that advances in LLM capabilities since the initial manual study suggest that automated assessment should closely approximate the results achieved through human review, while significantly reducing evaluation time and enabling larger-scale studies. Additionally, we expect that open-source models available through cost-effective providers may offer reasonable performance at a fraction of the cost of proprietary models.
 
\subsection{Related Work}
 
Prior research has examined vulnerabilities in LLM-generated code from various angles:
 
\citet{pearce2022asleep} demonstrated that GitHub Copilot generates functional but insecure code in certain scenarios, highlighting vulnerabilities in cryptographic usage and SQL injection risks. \citet{sandoval2023empirical} provided an empirical analysis of security vulnerabilities in code generated by various LLMs. \citet{Kolosnjaji2024} investigated how LLMs can be used to detect and evaluate code vulnerabilities.
 
However, to our knowledge, no prior work has conducted a systematic comparison of manual versus automated LLM-based assessment of code security. This represents a novel contribution: we directly compare the outcomes of human expert assessment with automated AI-driven evaluation, providing insights into the viability of scaling security audits through automation.
 
\section{Background: NIST Standards and Security Best Practices}
 
To properly contextualize our assessment framework, we provide detailed descriptions of the NIST guidelines used to evaluate ChatGPT responses.
 
\subsection{NIST Cybersecurity Framework}
 
The NIST Cybersecurity Framework provides a set of guidelines and best practices for organizations to manage and reduce cybersecurity risks \citep{nist2018framework}. The framework organizes practices into five core functions: Identify, Protect, Detect, Respond, and Recover. For the purpose of evaluating developer-LLM conversations, the most relevant function is \textit{Protect}, which encompasses practices for limiting or containing the impact of a potential cybersecurity event.
 
Key categories within the Protect function include:
\begin{itemize}
    \item \textbf{Access Control (PR.AC)}: Practices for controlling and limiting access to systems, including authentication, authorization, and privilege management.
    \item \textbf{Data Security (PR.DS)}: Practices for protecting data at rest and in transit, including encryption, data classification, and secure deletion.
    \item \textbf{Information Protection Processes (PR.IP)}: Practices for managing security updates, secure development, and secure configuration management.
    \item \textbf{Supply Chain Risk Management (PR.SY)}: Practices for managing risks from third-party code and dependencies.
\end{itemize}
 
\subsection{Specific Standards and Best Practices}
 
Our assessment codebook specifically evaluates ChatGPT responses against the following concrete standards:
 
TODO
 
\section{Methodology}
 
\subsection{Dataset and Prior Work}
 
Our dataset consists of 338 developer conversations between users and ChatGPT (GPT-3.5 model) from a publicly available Hugging Face dataset containing approximately 90,000 such conversations \citep{huggingface2023}. These 338 conversations were manually annotated by our team two semesters ago using a structured codebook that identified security and privacy-related discussions.
 
The original manual assessment process involved:
\begin{enumerate}
    \item Classification of conversations by topic (security, privacy, general development, etc.)
    \item Identification of relevant security/privacy questions
    \item Manual evaluation of ChatGPT's responses against NIST guidelines
    \item Documentation of compliance or violations
\end{enumerate}
 
\subsection{Data Ethics, Privacy, and Permissions}
 
\subsubsection{Use of Human Data and Anonymity}
 
The original Hugging Face dataset was created from public ChatGPT conversations, and we have verified that:
\begin{itemize}
    \item All user conversations in the dataset have been anonymized, with personal identifiers removed.
    \item The dataset does not contain sensitive personal information beyond the nature of the code/security questions being asked.
    \item The dataset is published under an open license permitting research use.
\end{itemize}
 
For our automated assessment work, we will not collect new human data, and we will not attempt to identify the original users in the Hugging Face dataset. Instead, we use only the anonymized conversations already published.
 
\subsubsection{Use of Previous Work Data}
 
Permission to use the previous project's manual annotations has been obtained from the team members involved in that project. The annotations themselves do not contain sensitive information—they consist of security/privacy classifications and compliance assessments. We confirm that using this prior work for comparison is appropriate under the research collaboration agreement.
 

 
\subsection{Experimental Design}
 
\subsubsection{Phase 1: Automated Assessment with Multiple LLMs}
 
We will assess the same 338 conversations using three categories of LLMs:

\begin{enumerate}
    \item \textbf{Proprietary, Expensive Models}:
    \begin{itemize}
        \item Claude (via Anthropic API)
        \item GPT-4 (via OpenAI API)
    \end{itemize}
    
    \item \textbf{Open-Source Models (via OpenRouter)}:
    \begin{itemize}
        \item Mistral 7B
        \item Llama 2 (7B and 13B variants)
        \item Other open-weight models as feasible
    \end{itemize}
    
    \item \textbf{Baseline Comparisons}:
    \begin{itemize}
        \item GPT-3.5 (the original model used in conversations)
    \end{itemize}
\end{enumerate}
 
For each conversation, we prompt the evaluating LLM with:
\begin{lstlisting}
You are a security expert evaluating code suggestions from 
ChatGPT against NIST Cybersecurity Framework standards.
 
Evaluate the following conversation against the criteria:
[Criteria listed]
 
CONVERSATION:
[User question and ChatGPT response]
 
Provide:
1. Overall compliance: COMPLIANT / NON-COMPLIANT / PARTIALLY COMPLIANT
2. Specific issues identified
3. Recommended improvements
4. Confidence score (0-100)
\end{lstlisting}
 
\subsubsection{Phase 2: Comparison Against Manual Annotations}
 
We will compare automated assessments against the manual annotations from the prior project using the following metrics:
 
\begin{itemize}
    \item \textbf{Agreement Rate}: Percentage of conversations where the automated assessment reaches the same compliance conclusion as manual review.
    \item \textbf{Precision and Recall}: Among conversations identified as non-compliant by manual review, what fraction are correctly identified by each automated method?
    \item \textbf{False Positive Rate}: How often do automated methods flag compliant code as non-compliant?
    \item \textbf{Confidence Distribution}: Do more confident LLM assessments correlate with higher accuracy? Use a Pearson Correlation Test
\end{itemize}
 
 
\subsection{Reproducibility and Code}
 
All code for prompting, data processing, and evaluation will be maintained in a public GitHub repository. 

Here is access to our old repository that performed the manual work: https://github.com/DishitaS123/HumanFac

New Repository: https://github.com/DishitaS123/DecodingGPT/tree/main

Do not quite have results yet, i was super swamped this week, but I will get started on results tomorrow.  

The anonymized dataset and manual annotations will be shared (subject to any applicable licensing restrictions).
 
\section{Preliminary Results and Findings}
 
\textit{[This section will be populated as experiments progress. The following is a placeholder describing what this section will contain.]}
 
At this stage, we are in the early phases of automated assessment. We have:
 
\begin{enumerate}
    \item Finalized the assessment prompt and rubric
    \item Set up API access to Claude and OpenRouter for cost-effective testing
    \item Begun running assessments on a subset of 50 conversations to validate the approach
\end{enumerate}
 
Preliminary observations:
\begin{itemize}
    \item Initial assessments suggest good agreement between Claude and manual annotations on clear-cut compliance cases
    \item The open-source models show some variability in their compliance assessments, which we are investigating
    \item Cost differences are substantial (Claude costs approximately 10x the cost of Llama 2 via OpenRouter)
\end{itemize}
 
We expect to complete the full assessment of all 338 conversations by end of next week, with final results presented in the conclusion.
 
 
\section{Conclusion}
 
\textit{[This section will be completed when results are finalized. The following is a placeholder.]}
 
% This project addresses a gap in the literature by providing the first systematic comparison of automated versus manual assessment of LLM-generated code security. By evaluating multiple modern LLMs across 338 developer conversations, we aim to establish whether AI-driven security assessment can reliably replicate expert human judgment while dramatically reducing cost and time.
 
% The implications of this work extend beyond academic interest: as LLMs become more prevalent in professional software development, organizations need scalable methods for auditing the security of suggestions provided to developers. This project contributes a methodology and empirical evidence toward that goal.
 
\begin{thebibliography}{99}
 
\bibitem{nist2018framework} NIST. (2018). \textit{Cybersecurity Framework Version 1.1}. National Institute of Standards and Technology. Retrieved from https://www.nist.gov/cyberframework
 
\bibitem{pearce2022asleep} Pearce, H., Tan, B., Rahman, B., Dhananjay, A., \& Paxson, V. (2022). Asleep at the Keyboard? Assessing the Security of GitHub Copilot's Code Contributions. In \textit{2022 IEEE Symposium on Security and Privacy (SP)} (pp. 754-774). IEEE.
 
\bibitem{sandoval2023empirical} Sandoval, G., Kouton, H., Chow, J., \& Wagner, D. (2023). Empirical Analysis of Code Injection Attacks in GitHub Copilot Suggestions. In \textit{Proceedings of the 32nd USENIX Security Symposium}.
 
\bibitem{Kolosnjaji2024} Kolosnjaji, B., Eckert, C., \& Zarras, A. (2024). Towards Understanding and Improving Language Model-Based Vulnerability Detection. In \textit{Proceedings of the 29th European Symposium on Research in Computer Security (ESORICS)}.
 
\bibitem{huggingface2023} Hugging Face. (2023). Conversations with ChatGPT Dataset. Retrieved from https://huggingface.co/datasets/
 
\end{thebibliography}
 
\appendix
 
% \section{NIST Framework Details}
 
% \subsection{Full Criteria Reference}
 
% \textit{[Detailed expanded list of all NIST criteria being evaluated. This can be expanded as the project progresses.]}
 
% \section{Sample Assessment Prompts}
 
% This section documents the exact prompts used to query each LLM for assessment. Final versions will be included here.
 
% \section{Preliminary Results Tables}
 
% \textit{[Tables and figures comparing initial assessments across models will be added as results accumulate.]}
 
\end{document}